% Vietnamese translation for OI Public Library
% Source-File: IOI中国国家候选队论文/2007/day2/8.陈瑜希《多角度思考 创造性思维》.ppt
\documentclass[11pt]{article}
\usepackage{oipl-vn}

\newcommand{\Oh}{\mathcal{O}}
\newcommand{\dist}{\operatorname{dist}}

\title{Bản ghi chú trình chiếu: Suy nghĩ từ nhiều góc độ, tư duy sáng tạo}
\author{Chen Yuxi}
\date{2007}

\begin{document}
\maketitle

\origtitle{Suy nghĩ từ nhiều góc độ, tư duy sáng tạo}
\sourcepath{IOI China candidate team papers / 2007 / day 2 / Chen Yuxi.ppt}

\begin{abstract}
Các slide dùng một số bài quy hoạch động trên cây để minh họa cách đổi góc
nhìn khi trạng thái tự nhiên chưa đủ. Mục tiêu không chỉ là trình bày công
thức, mà còn cho thấy vì sao phải chọn điểm phân nhánh, nhà máy tổ tiên, dãy
nhãn tổ tiên hoặc trọng tâm cây làm đối tượng trung tâm.
\end{abstract}

\paragraph{Từ khóa.}
Cấu trúc cây; quy hoạch động; truyền thông tin; đổi gốc; tư duy sáng tạo.

\section{Khung chung của DP trên cây}

Khi bài toán yêu cầu tối ưu trên cây, cách làm thường gồm ba bước:
\begin{enumerate}
  \item xác lập trạng thái: cần lưu thông tin nào của cây con, có phải thêm
  chiều hay không;
  \item chuyển trạng thái: gộp các con, truyền thông tin từ cha xuống con,
  hoặc phối hợp cả hai hướng;
  \item cài đặt: DFS có nhớ, BFS theo độ sâu, hoặc hai lần duyệt.
\end{enumerate}

Khác với DP tuyến tính chỉ có chiều tiến hoặc lùi, DP trên cây có hai hướng
tự nhiên: từ lá lên gốc và từ gốc xuống lá. Nhiều bài cần cả hai. Điểm khó
thường không nằm ở việc viết DFS, mà ở việc nhìn đúng vai trò của một đỉnh
trong lời giải.

\section{Ví dụ 1: Tìm Chris}

\subsection{Mô hình}

Thành phố là một cây có trọng số. Cần chọn ba đỉnh \(X,Y,Z\) sao cho
\[
  \dist(X,Y)\le \dist(X,Z)
\]
và
\[
  \dist(X,Y)+\dist(Y,Z)
\]
lớn nhất. Ở đây \(X\) là nhà Chris, \(Y\) là nơi bố mẹ đi tìm trước, còn
\(Z\) là nơi đi tìm sau trong trường hợp xấu nhất.

Nếu cố định \(Y\), ta cần tìm hai đỉnh xa nhất liên quan tới \(Y\). Làm vậy
cho mọi \(Y\) mất \(\Oh(n^2)\), không chạy được với \(n\) lớn. Slide đổi góc
nhìn: thay vì liệt kê nơi đến trước, liệt kê \term{điểm phân nhánh} của hai
đường đi.

\subsection{Điểm phân nhánh}

Giả sử \(a\) là điểm mà hai đường \(Y\to X\) và \(Y\to Z\) tách nhau. Gốc hóa
cây tại \(a\). Ba đỉnh \(X,Y,Z\) nằm ở ba nhánh khác nhau đi ra từ \(a\);
trường hợp \(Y=a\) có thể xem như một nhánh rỗng độ dài \(0\). Khi đó
\[
\begin{aligned}
  \dist(X,Y)+\dist(Y,Z)
  &=\dist(X,a)+2\dist(Y,a)+\dist(Z,a).
\end{aligned}
\]
Điều kiện \(\dist(X,Y)\le\dist(X,Z)\) tương đương
\(\dist(Y,a)\le\dist(Z,a)\). Vì vậy, nếu ba khoảng cách lớn nhất từ \(a\) tới
ba nhánh khác nhau là \(d_1\ge d_2\ge d_3\), giá trị tốt nhất tại \(a\) là
\[
  d_1+2d_2+d_3.
\]

\subsection{Hai lần duyệt}

Gốc hóa cây tại một đỉnh tùy ý. Lần duyệt thứ nhất đi từ lá lên gốc để tính
các khoảng cách tốt nhất nằm trong cây con. Lần duyệt thứ hai đi từ gốc xuống
lá để truyền phần tốt nhất nằm ngoài cây con. Ở mỗi đỉnh, lưu vài giá trị lớn
nhất và lưu mỗi giá trị đến từ cạnh kề nào.

Khi truyền từ \(u\) sang con \(v\), nếu giá trị lớn nhất của \(u\) đến từ
\(v\) thì dùng giá trị lớn thứ hai; nếu không, dùng giá trị lớn nhất. Cộng
thêm độ dài cạnh \(uv\) để được ứng viên nhìn từ \(v\). Sau hai lần duyệt,
mỗi đỉnh có đủ các ứng viên thuộc các nhánh khác nhau, và ta cập nhật đáp án
bằng công thức \(d_1+2d_2+d_3\). Cả thời gian và bộ nhớ là \(\Oh(n)\).

Ý tưởng đáng học là: một đại lượng khó khi nhìn từ đầu mút lại trở nên đơn
giản khi nhìn từ điểm phân nhánh.

\section{Ví dụ 2: Nhà máy gỗ ở Byteland}

\subsection{Mô hình}

Các làng nằm trên hệ thống sông không phân nhánh theo chiều xuôi dòng; gốc là
Bytetown. Cần xây thêm \(k\) nhà máy gỗ ở các làng để mỗi làng đưa gỗ tới nhà
máy đầu tiên gặp trên đường về gốc, sao cho tổng chi phí vận chuyển nhỏ nhất.

Trạng thái chỉ gồm đỉnh hiện tại và số nhà máy trong cây con là chưa đủ, vì
chi phí của một làng còn phụ thuộc vào nhà máy tổ tiên gần nhất ở phía trên.
Vì vậy slide thêm một chiều trạng thái cho ``nhà máy tổ tiên đang được truyền
xuống''.

\subsection{Trạng thái và chuyển}

Đặt
\[
  f(u,t,q)
\]
là chi phí nhỏ nhất trong cây con gốc \(u\), khi đặt \(q\) nhà máy trong cây
con này, và \(t\) là một tổ tiên đã có nhà máy. Có hai trường hợp.

Nếu đặt nhà máy tại \(u\), gỗ của \(u\) không tốn chi phí, còn các con dùng
\(u\) làm nhà máy tổ tiên. Một nhà máy đã dùng ở \(u\); số nhà máy còn lại
được phân phối giữa các con, giống ba lô.

Nếu không đặt nhà máy tại \(u\), gỗ của \(u\) phải đi tới \(t\), sinh chi phí
\(w_u\dist(u,t)\), và các con tiếp tục dùng \(t\). Nếu \(u\) có các con
\(\operatorname{son}_i\), quá trình gộp con dùng các mảng tạm:
\[
  h(i,s)=\min_{0\le r\le s}
  \{h(i-1,s-r)+f(\operatorname{son}_i,\text{tổ tiên mới},r)\}.
\]
Sau khi gộp đủ con, lấy nhỏ nhất giữa hai trường hợp đặt và không đặt nhà máy
tại \(u\).

Bài học ở đây là một chiều trạng thái tưởng như ``bên ngoài cây con'' lại là
thứ làm cây con trở nên độc lập, từ đó chuyển trạng thái rõ ràng hơn.

\section{Ví dụ 3: Tính phí mạng theo cặp}

\subsection{Biến đổi chi phí}

Mạng người dùng là cây nhị phân đầy đủ, lá là người dùng. Mỗi người chọn
phương thức \(A\) hoặc \(B\). Với mỗi cặp người dùng, hệ số tính phí phụ thuộc
vào lựa chọn của hai đầu mút và vào so sánh số người chọn \(A\), \(B\) trong
miền của LCA của cặp đó.

Bảng hệ số có dạng đặc biệt. Khi trong miền LCA có \(n_A<n_B\), mỗi đầu mút
chọn \(A\) đóng một lần lưu lượng; khi \(n_A\ge n_B\), mỗi đầu mút chọn \(B\)
đóng một lần lưu lượng. Nhờ vậy, phí theo cặp có thể được chuyển thành phí
gán cho từng người dùng theo từng tổ tiên. Đây là bước sáng tạo chính trước
khi đặt DP.

\subsection{Mã hóa trạng thái}

Với một đỉnh \(i\) của cây nhị phân, trạng thái cần biết:
\begin{itemize}
  \item trong miền của \(i\) có bao nhiêu lá chọn \(A\);
  \item với từng tổ tiên của \(i\), điều kiện \(n_A<n_B\) hay \(n_A\ge n_B\)
  đang đúng.
\end{itemize}
Gọi dãy nhãn tổ tiên này là một số nhị phân \(k\). Đặt
\[
  f(i,j,k)
\]
là chi phí nhỏ nhất trong miền của \(i\), khi có đúng \(j\) người chọn \(A\)
và dãy nhãn tổ tiên là \(k\).

Khi gộp hai con \(L,R\), phân phối \(j\) người chọn \(A\): nếu con trái có
\(u\) người chọn \(A\) thì con phải có \(j-u\). Từ \(j\) và kích thước miền
của \(i\), ta biết nhãn mới của \(i\), rồi đẩy nhãn này vào mã \(k\) khi đi
xuống hai con:
\[
  f(i,j,k)=
  \min_u\{f(L,u,\operatorname{push}(k,s))
        +f(R,j-u,\operatorname{push}(k,s))+\text{phí qua }i\}.
\]
Ở lá, mã \(k\) cho biết lá phải đóng phí qua các tổ tiên theo phương thức
nào, rồi cộng thêm chi phí đổi phương thức nếu khác đăng ký ban đầu.

Số trạng thái nhìn thô rất lớn, nhưng theo tầng, mỗi đỉnh chỉ có
\(\Oh(2^{N-i})\) giá trị của \(j\) và \(\Oh(2^i)\) mã tổ tiên. Tổng hợp lại,
không gian có thể giữ ở \(\Oh(m^2)\) và thời gian ở \(\Oh(m^2N)\) với
\(m=2^N\).

\section{Ví dụ 4: Hang động và đoán kho báu}

Trong một cây, A giấu kho báu ở một phòng. Mỗi lần B đoán một phòng; nếu sai,
A chỉ cho biết hành lang đầu tiên trên đường từ phòng vừa đoán tới kho báu.
Cần số lần đoán ít nhất trong trường hợp xấu nhất.

Nếu đoán một đỉnh \(u\), phản hồi của A chia cây thành các thành phần sau khi
xóa \(u\): kho báu nằm trong một thành phần cụ thể, hoặc chính là \(u\). Vì
vậy một chiến lược tốt cố gắng chọn đỉnh làm cân bằng kích thước trường hợp
xấu nhất. Góc nhìn tự nhiên là trọng tâm cây: chọn đỉnh sao cho thành phần
lớn nhất sau khi xóa nó nhỏ nhất. Sau mỗi phản hồi, bài toán thu hẹp vào một
cây con, và tiếp tục dùng cùng ý tưởng.

Slide dùng ví dụ này để nhấn mạnh rằng không phải mọi bài cây đều nên đặt
trạng thái theo ``cây con gốc tại \(u\)''. Có lúc đối tượng cần nhìn là đỉnh
chia tách thông tin tốt nhất trong quá trình hỏi đáp.

\section{Kết luận}

Quy hoạch động trên cây không chỉ là thay chỉ số tuyến tính bằng chỉ số đỉnh.
Trong bốn ví dụ, lời giải xuất hiện khi ta đổi góc nhìn: điểm phân nhánh thay
cho đầu mút, nhà máy tổ tiên thay cho cây con độc lập, dãy nhãn tổ tiên thay
cho một biến cục bộ, và trọng tâm cây thay cho gốc cố định. Khi trạng thái tự
nhiên quá thiếu hoặc quá lớn, cần hỏi lại: thông tin thật sự được truyền qua
cạnh nào, và vai trò nào của đỉnh làm bài toán đơn giản nhất?

\end{document}
